% ============================================================
%  research-report — shared preamble
%  Compile with pdflatex (portable) or lualatex/xelatex.
%  Only packages present in a minimal TeX Live install are
%  required; optional upgrades are guarded with \IfFileExists.
% ============================================================

% ---- page geometry -----------------------------------------
\usepackage[letterpaper,margin=1.15in,footskip=0.5in]{geometry}
% Text width is ~5.5in — figures are built at that width.

% ---- encoding and fonts ------------------------------------
\usepackage[T1]{fontenc}
\usepackage[utf8]{inputenc}
\IfFileExists{charter.sty}{\usepackage{charter}}{}   % readable serif body
\IfFileExists{newtxmath.sty}{\usepackage[charter]{newtxmath}}{}
\usepackage[scaled=0.86]{helvet}                     % sans for headings
\usepackage{microtype}

% ---- colour ------------------------------------------------
\usepackage[dvipsnames]{xcolor}
\definecolor{rrInk}    {HTML}{1A1A1A}   % body text
\definecolor{rrMuted}  {HTML}{5A5F66}   % metadata, captions
\definecolor{rrAccent} {HTML}{0F4C81}   % headings, links
\definecolor{rrRule}   {HTML}{C9CED6}   % hairlines
\definecolor{rrBoxBG}  {HTML}{F4F6F8}   % callout background
\color{rrInk}

% ---- spacing and lists -------------------------------------
\usepackage{setspace}
\setstretch{1.10}
\usepackage{enumitem}
\setlist{noitemsep,topsep=4pt,leftmargin=1.4em}
\usepackage{parskip}

% ---- headings ----------------------------------------------
\usepackage{titlesec}
\titleformat{\section}
  {\normalfont\sffamily\bfseries\large\color{rrAccent}}{\thesection}{0.7em}{}
\titleformat{\subsection}
  {\normalfont\sffamily\bfseries\normalsize\color{rrInk}}{\thesubsection}{0.6em}{}
\titleformat{\subsubsection}
  {\normalfont\sffamily\itshape\normalsize\color{rrInk}}{}{0em}{}
\titlespacing*{\section}   {0pt}{18pt plus 4pt minus 2pt}{6pt}
\titlespacing*{\subsection}{0pt}{12pt plus 3pt minus 2pt}{4pt}

% ---- headers and footers -----------------------------------
\usepackage{fancyhdr}
\pagestyle{fancy}
\fancyhf{}
\renewcommand{\headrulewidth}{0.4pt}
\renewcommand{\headrule}{\hbox to\headwidth{\color{rrRule}\leaders\hrule height \headrulewidth\hfill}}
\fancyhead[L]{\footnotesize\sffamily\color{rrMuted}\rrShortTitle}
\fancyhead[R]{\footnotesize\sffamily\color{rrMuted}\rrStatus\ \textbullet\ \rrDate}
\fancyfoot[C]{\footnotesize\sffamily\color{rrMuted}\thepage}

% ---- tables ------------------------------------------------
\usepackage{booktabs}
\usepackage{array}
\usepackage{longtable}
\renewcommand{\arraystretch}{1.15}
\newcolumntype{L}[1]{>{\raggedright\arraybackslash}p{#1}}
\newcolumntype{C}[1]{>{\centering\arraybackslash}p{#1}}
\newcolumntype{R}[1]{>{\raggedleft\arraybackslash}p{#1}}
% numeric column aligned on a fixed decimal count: use \num{}-free plain digits

% ---- figures and captions ----------------------------------
\usepackage{graphicx}
\usepackage{float}
\usepackage[font={small,color=rrMuted},labelfont={bf,color=rrInk},
            labelsep=period,skip=6pt,width=0.95\textwidth]{caption}

% ---- callout boxes -----------------------------------------
\usepackage[most]{tcolorbox}
\newtcolorbox{rrsummary}[1][]{
  colback=rrBoxBG, colframe=rrAccent, boxrule=0pt, leftrule=2.5pt,
  arc=0pt, outer arc=0pt, left=10pt, right=10pt, top=8pt, bottom=8pt,
  breakable, #1}
\newtcolorbox{rrnote}[1][]{
  colback=white, colframe=rrRule, boxrule=0.6pt,
  arc=1pt, left=8pt, right=8pt, top=6pt, bottom=6pt,
  breakable, #1}

% ---- bibliography ------------------------------------------
\usepackage[round,authoryear]{natbib}
\bibliographystyle{plainnat}
\setlength{\bibsep}{3pt}

% ---- links -------------------------------------------------
\usepackage[hidelinks,breaklinks=true]{hyperref}
\hypersetup{colorlinks=true, linkcolor=rrAccent, citecolor=rrAccent,
            urlcolor=rrAccent, pdfborder={0 0 0}}
\usepackage{url}
\urlstyle{same}

% ============================================================
%  Report metadata — set these in report.tex
% ============================================================
\newcommand{\rrTitle}{Untitled report}
\newcommand{\rrShortTitle}{Untitled}
\newcommand{\rrAuthors}{Unknown}
\newcommand{\rrDate}{\today}
\newcommand{\rrStatus}{Draft}
\newcommand{\rrSlug}{untitled}
\newcommand{\rrRelated}{---}
\newcommand{\rrArtifacts}{---}

\newcommand{\reporttitle}[2]{\renewcommand{\rrTitle}{#1}\renewcommand{\rrShortTitle}{#2}}
\newcommand{\reportauthors}[1]{\renewcommand{\rrAuthors}{#1}}
\newcommand{\reportdate}[1]{\renewcommand{\rrDate}{#1}}
\newcommand{\reportstatus}[1]{\renewcommand{\rrStatus}{#1}}
\newcommand{\reportslug}[1]{\renewcommand{\rrSlug}{#1}}
\newcommand{\reportrelated}[1]{\renewcommand{\rrRelated}{#1}}
\newcommand{\reportartifacts}[1]{\renewcommand{\rrArtifacts}{#1}}

% ---- title block -------------------------------------------
\newcommand{\makereporthead}{%
  \thispagestyle{fancy}
  \begingroup
  \raggedright
  {\sffamily\bfseries\LARGE\color{rrAccent}\rrTitle\par}
  \vspace{6pt}
  {\sffamily\small\color{rrMuted}\rrAuthors\ \textbullet\ \rrDate\ \textbullet\ Status: \textbf{\rrStatus}\par}
  \vspace{8pt}
  {\color{rrRule}\hrule height 0.8pt}
  \vspace{8pt}
  {\footnotesize\sffamily\color{rrMuted}
   \begin{tabular}{@{}p{0.14\textwidth}p{0.80\textwidth}@{}}
     Slug      & \texttt{\rrSlug} \\
     Related   & \rrRelated \\
     Artifacts & \rrArtifacts \\
   \end{tabular}\par}
  \vspace{4pt}
  {\color{rrRule}\hrule height 0.4pt}
  \vspace{12pt}
  \endgroup
}

% ---- confidence tags ---------------------------------------
% Usage: \confidence{suggestive} inline in the Discussion.
\newcommand{\confidence}[1]{%
  \textsf{\footnotesize\color{rrAccent}[\,confidence: #1\,]}}

% ---- provenance pointer ------------------------------------
% Usage: \source{analysis/pooling_recall.py}
% \detokenize lets you write underscores, tildes etc. without escaping them.
\newcommand{\source}[1]{%
  \textsf{\footnotesize\color{rrMuted}(\texttt{\detokenize{#1}})}}

% ---- inline path / identifier ------------------------------
% Same escaping-free behaviour for commands, configs, run IDs.
\newcommand{\pathname}[1]{\texttt{\detokenize{#1}}}

% ---- standard figure include -------------------------------
% \rrfigure{path}{caption}{label}
\newcommand{\rrfigure}[3]{%
  \begin{figure}[htbp]
    \centering
    \includegraphics{#1}
    \caption{#2}
    \label{#3}
  \end{figure}}
